\chapter{Coefficienti binomiali}\label{cap:binomiali}

Quante sono le parti con esattamente $k$ elementi di un insieme con $n$ elementi? Ad esempio, in quanti modi si scelgono
$3$ persone su $10$ per formare una commissione? La risposta è un numero che compare ovunque in matematica: il
\textbf{coefficiente binomiale} $\binom{n}{k}$.

% ──────────────────────────────────────────────────────────────
\section{Definizione e prime proprietà}\label{sec:def-binomiali}

\dfn[binomiale]{$k$-parti e coefficiente binomiale}{
    Sia $s$ un insieme e $k \in \mbN$. Indichiamo con $\mcP_k(s) := \{t \in \mcP(s) \mid |t| = k\}$ l'insieme delle parti di $s$ con
    esattamente $k$ elementi (le \textbf{$k$-parti} di $s$). Per $n, k \in \mbN$ il \textbf{coefficiente binomiale} «$n$ su $k$» è
    \[ \binom{n}{k} := |\mcP_k(I_n)|, \]
    il numero di parti con $k$ elementi di $\{1, \dots, n\}$.
}

\begin{esempio}{Primi valori}
    Per ogni $n$: $\binom{n}{0} = 1$ (l'unica $0$-parte è $\varnothing$), $\binom{n}{n} = 1$ (l'unica $n$-parte di $I_n$ è $I_n$), $\binom{n}{1} = n$
    (i singleton). Se $k > n$, $\binom{n}{k} = 0$, perché una parte di $I_n$ ha al più $n$ elementi. Le $2$-parti di $I_4$ sono
    $\{1, 2\}, \{1, 3\}, \{1, 4\}, \{2, 3\}, \{2, 4\}, \{3, 4\}$, quindi $\binom{4}{2} = 6$.
\end{esempio}

La definizione usa $I_n$, ma vale per qualunque insieme con $n$ elementi: contano solo quanti sono, non come si chiamano.

\begin{concetto}
\begin{Proposizione}{Le $k$-parti di un insieme qualunque}{k-parti}
    \begin{prerequisiti}
        \dip{Definizione}{def:binomiale}
        \dip{Definizione}{def:immagine-antiimmagine}
        \dip{Proposizione}{prop:ridotta}
    \end{prerequisiti}
    Se $|s| = n$, allora $|\mcP_k(s)| = \binom{n}{k}$ per ogni $k \in \mbN$.
\end{Proposizione}
\begin{dimostrazione}
    Sia $f \colon s \to I_n$ biettiva. Per ogni $t \subseteq s$, la restrizione di $f$ a $t$ è una biezione tra $t$ e $\overrightarrow{f}(t)$, quindi
    $|\overrightarrow{f}(t)| = |t|$. Allora $t \mapsto \overrightarrow{f}(t)$ manda $\mcP_k(s)$ in $\mcP_k(I_n)$, ed è biettiva: la sua inversa è
    $u \mapsto \overrightarrow{f^{-1}}(u)$. Quindi $|\mcP_k(s)| = |\mcP_k(I_n)| = \binom{n}{k}$.
\end{dimostrazione}
\end{concetto}

\begin{spiegazione}
    Una biezione è un cambio di nomi: rinominando gli elementi di $s$ come $1, \dots, n$, le $k$-parti di $s$ diventano le $k$-parti di
    $I_n$, una per una. Ad esempio le $2$-parti di $\{a, b, c, d\}$ sono $6$ come quelle di $I_4$.
\end{spiegazione}

\begin{concetto}
\begin{Teorema}{La somma dei coefficienti binomiali}{somma-binomiali}
    \begin{prerequisiti}
        \dip{Definizione}{def:binomiale}
        \dip{Definizione}{def:partizione}
        \dip{Lemma}{lem:somma-prodotto}
        \dip{Corollario}{cor:cardinalita-parti}
    \end{prerequisiti}
    Per ogni $n \in \mbN$: \quad $\displaystyle\sum_{k = 0}^{n} \binom{n}{k} = \binom{n}{0} + \binom{n}{1} + \dots + \binom{n}{n} = 2^n$.
\end{Teorema}
\begin{dimostrazione}
    Ogni parte di $I_n$ ha un numero di elementi $k$ compreso tra $0$ e $n$, e uno solo; inoltre $\mcP_k(I_n) \neq \varnothing$ per $0 \leq k \leq n$.
    Quindi $\{\mcP_k(I_n) \mid 0 \leq k \leq n\}$ è una partizione di $\mcP(I_n)$ in $n + 1$ blocchi, e per il Lemma~\ref{lem:somma-prodotto}
    \[ \sum_{k = 0}^{n} |\mcP_k(I_n)| = |\mcP(I_n)| = 2^n \]
    (Corollario~\ref{cor:cardinalita-parti}).
\end{dimostrazione}
\end{concetto}

\begin{spiegazione}
    Si contano le parti di $I_n$ \textbf{raggruppandole per numero di elementi}. Per $n = 3$ i gruppi sono le righe del diagramma
    di Hasse di $\mcP(\{a, b, c\})$:
    \begin{center}
    \begin{tikzpicture}[hasse]
        \node[hv] (v1) at (0,0) {$\varnothing$};
        \node[hv evid] (a) at (-1.5,1.1) {$\{a\}$}; \node[hv evid] (b) at (0,1.1) {$\{b\}$}; \node[hv evid] (c) at (1.5,1.1) {$\{c\}$};
        \node[hv evid2] (ab) at (-1.5,2.2) {$\{a, b\}$}; \node[hv evid2] (ac) at (0,2.2) {$\{a, c\}$}; \node[hv evid2] (bc) at (1.5,2.2) {$\{b, c\}$};
        \node[hv] (abc) at (0,3.3) {$\{a, b, c\}$};
        \draw[hlato] (v1) -- (a); \draw[hlato] (v1) -- (b); \draw[hlato] (v1) -- (c);
        \draw[hlato] (a) -- (ab); \draw[hlato] (a) -- (ac); \draw[hlato] (b) -- (ab); \draw[hlato] (b) -- (bc);
        \draw[hlato] (c) -- (ac); \draw[hlato] (c) -- (bc);
        \draw[hlato] (ab) -- (abc); \draw[hlato] (ac) -- (abc); \draw[hlato] (bc) -- (abc);
        \foreach \y/\k/\v in {0/0/1, 1.1/1/3, 2.2/2/3, 3.3/3/1}
            \node[font=\small, anchor=west] at (2.6,\y) {$\binom{3}{\k} = \v$};
        \node[font=\small, anchor=west] at (4.6,1.65) {$1 + 3 + 3 + 1 = 8 = 2^3$};
    \end{tikzpicture}
    \end{center}
\end{spiegazione}

\begin{concetto}
\begin{Teorema}{Simmetria dei coefficienti binomiali}{simmetria-binomiali}
    \begin{prerequisiti}
        \dip{Definizione}{def:binomiale}
        \dip{Teorema}{teo:composte-caratteristiche}
        \dip{Teorema}{teo:inclusione-esclusione}
        \dip{Proposizione}{prop:doppia-neg-ins}
    \end{prerequisiti}
    Per ogni $n, k \in \mbN$ con $k \leq n$: \quad $\displaystyle\binom{n}{k} = \binom{n}{n - k}$.
\end{Teorema}
\begin{dimostrazione}
    La funzione $x \in \mcP(I_n) \mapsto I_n \setminus x \in \mcP(I_n)$ è biettiva: applicata due volte dà $I_n \setminus (I_n \setminus x) = x$ (doppia
    negazione insiemistica), quindi è l'inversa di sé stessa. Se $|x| = k$, allora $|I_n \setminus x| = n - k$, perché $I_n$ è unione disgiunta di
    $x$ e $I_n \setminus x$. Quindi la funzione manda $\mcP_k(I_n)$ in $\mcP_{n - k}(I_n)$, e la sua restrizione è una biezione tra i due
    (con inversa la restrizione $\mcP_{n - k}(I_n) \to \mcP_k(I_n)$ della stessa funzione). Dunque $|\mcP_k(I_n)| = |\mcP_{n - k}(I_n)|$.
\end{dimostrazione}
\end{concetto}

\begin{spiegazione}
    Scegliere i $k$ elementi \textbf{da prendere} equivale a scegliere gli $n - k$ elementi \textbf{da lasciare}. Scegliere $3$ persone su
    $10$ per una commissione è come scegliere le $7$ che restano fuori: $\binom{10}{3} = \binom{10}{7}$. Nel diagramma di Hasse qui sopra,
    il complementare scambia la riga $k$ con la riga $n - k$: è il diagramma capovolto.
\end{spiegazione}

% ──────────────────────────────────────────────────────────────
\section{La formula ricorsiva e il triangolo di Tartaglia}\label{sec:tartaglia}

\begin{concetto}
\begin{Teorema}{Formula ricorsiva dei coefficienti binomiali}{stifel}
    \begin{prerequisiti}
        \dip{Definizione}{def:binomiale}
        \dip{Lemma}{lem:somma-prodotto}
        \dip{Lemma}{lem:togli-elemento}
    \end{prerequisiti}
    Per ogni $n, k \in \mbN$ con $k \leq n$: \quad $\displaystyle\binom{n + 1}{k + 1} = \binom{n}{k + 1} + \binom{n}{k}$.
\end{Teorema}
\begin{dimostrazione}
    Dividiamo le $(k + 1)$-parti di $I_{n + 1}$ in due gruppi, secondo se contengono o no l'ultimo elemento $n + 1$:
    \[ a = \{x \in \mcP_{k + 1}(I_{n + 1}) \mid n + 1 \notin x\}, \qquad b = \{x \in \mcP_{k + 1}(I_{n + 1}) \mid n + 1 \in x\}. \]
    $a$ e $b$ sono disgiunti e la loro unione è $\mcP_{k + 1}(I_{n + 1})$, quindi $\binom{n + 1}{k + 1} = |a| + |b|$ (Lemma~\ref{lem:somma-prodotto}).
    \begin{caso}{$|a| = \binom{n}{k + 1}$}
        Le parti di $I_{n + 1}$ che non contengono $n + 1$ sono le parti di $I_n$: quindi $a = \mcP_{k + 1}(I_n)$.
    \end{caso}
    \begin{caso}{$|b| = \binom{n}{k}$}
        $x \in b \mapsto x \setminus \{n + 1\} \in \mcP_k(I_n)$ è biettiva (Lemma~\ref{lem:togli-elemento}): la sua inversa aggiunge $n + 1$, cioè
        $y \mapsto y \cup \{n + 1\}$.
    \end{caso}
\end{dimostrazione}
\end{concetto}

\begin{spiegazione}
    Per scegliere $k + 1$ persone tra $n + 1$, guardiamo una persona particolare, l'ultima: o \textbf{non} viene scelta, e allora le
    $k + 1$ si scelgono tra le altre $n$; oppure viene scelta, e allora ne mancano $k$ da scegliere tra le altre $n$. Ad esempio
    $\binom{4}{2} = \binom{3}{2} + \binom{3}{1} = 3 + 3 = 6$: le $2$-parti di $I_4$ senza il $4$ sono $\{1,2\}, \{1,3\}, \{2,3\}$; quelle con il $4$ sono
    $\{1,4\}, \{2,4\}, \{3,4\}$.
\end{spiegazione}

La formula ricorsiva permette di calcolare tutti i coefficienti binomiali uno dopo l'altro, disponendoli nel
\textbf{triangolo di Tartaglia} (o di Pascal): la riga $n$ contiene $\binom{n}{0}, \binom{n}{1}, \dots, \binom{n}{n}$, i bordi valgono $1$ e
\textbf{ogni numero interno è la somma dei due numeri subito sopra di lui}.

\begin{concetto}
\begin{center}
\begin{tikzpicture}[x=0.62cm, y=0.62cm]
    \foreach \n in {0,...,6} {
        \foreach \k in {0,...,\n} {
            \pgfmathtruncatemacro{\v}{round(factorial(\n)/(factorial(\k)*factorial(\n-\k)))}
            \node[circle, minimum size=7mm, inner sep=0pt, font=\small, fill=primary!7, draw=primary!40] (t\n-\k) at (2*\k-\n, -1.6*\n) {$\v$};
        }
        \node[font=\small, anchor=east] at (-8.2, -1.6*\n) {$n = \n$};
    }
    \node[circle, minimum size=7mm, inner sep=0pt, font=\small, fill=orange!30, draw=orange!80!black] at (t3-1) {$3$};
    \node[circle, minimum size=7mm, inner sep=0pt, font=\small, fill=orange!30, draw=orange!80!black] at (t3-2) {$3$};
    \node[circle, minimum size=7mm, inner sep=0pt, font=\small, fill=seagreen!30, draw=seagreen!70!black] at (t4-2) {$6$};
    \draw[-{Stealth}, seagreen!70!black, line width=0.8pt] (t3-1) -- (t4-2);
    \draw[-{Stealth}, seagreen!70!black, line width=0.8pt] (t3-2) -- (t4-2);
    \node[font=\small, anchor=west] at (7.4,-6.4) {$\binom{4}{2} = \binom{3}{1} + \binom{3}{2} = 3 + 3$};
\end{tikzpicture}
\end{center}
Nella riga $n$ si leggono anche la simmetria $\binom{n}{k} = \binom{n}{n - k}$ (la riga è palindroma) e la somma $2^n$
(ad esempio $1 + 4 + 6 + 4 + 1 = 16 = 2^4$).
\end{concetto}

% ──────────────────────────────────────────────────────────────
\section{La formula chiusa}\label{sec:formula-chiusa}

\begin{concetto}
\begin{Teorema}{Formula chiusa dei coefficienti binomiali}{formula-binomiali}
    \begin{prerequisiti}
        \dip{Teorema}{teo:stifel}
        \dip{Definizione}{def:fattoriale}
        \dip{Corollario}{cor:induzione-uso}
    \end{prerequisiti}
    Per ogni $n, k \in \mbN$ con $k \leq n$: \quad $\displaystyle\binom{n}{k} = \frac{n!}{k!\,(n - k)!}$.
\end{Teorema}
\begin{dimostrazione}
    Per induzione su $n$, dimostrando la formula per tutti i $k \leq n$ contemporaneamente.
    \textbf{Base} $n = 0$: $\binom{0}{0} = 1 = \frac{0!}{0!\,0!}$. \textbf{Passo}: supponiamo la formula vera per $n$ e ogni $k \leq n$, e sia $k \leq n + 1$.
    \begin{caso}{$k = 0$ oppure $k = n + 1$}
        $\binom{n + 1}{0} = 1 = \frac{(n + 1)!}{0!\,(n + 1)!}$ e $\binom{n + 1}{n + 1} = 1 = \frac{(n + 1)!}{(n + 1)!\,0!}$.
    \end{caso}
    \begin{caso}{$1 \leq k \leq n$}
        Scriviamo $k = j + 1$ con $0 \leq j \leq n - 1$. Per la formula ricorsiva e l'ipotesi induttiva (applicata a $j + 1 \leq n$ e a $j \leq n$):
        \[ \binom{n + 1}{j + 1} = \binom{n}{j + 1} + \binom{n}{j} = \frac{n!}{(j + 1)!\,(n - j - 1)!} + \frac{n!}{j!\,(n - j)!}
           = \frac{n!\,\big((n - j) + (j + 1)\big)}{(j + 1)!\,(n - j)!} = \frac{(n + 1)!}{(j + 1)!\,(n - j)!}, \]
        dove abbiamo portato le due frazioni allo stesso denominatore $(j + 1)!\,(n - j)!$, moltiplicando la prima per $n - j$ e la
        seconda per $j + 1$. Poiché $j + 1 = k$ e $n - j = (n + 1) - k$, è la formula per $n + 1$ e $k$.
    \end{caso}
\end{dimostrazione}
\end{concetto}

\begin{spiegazione}
    Ecco un altro modo di capire la formula, che usa le funzioni iniettive (Teorema~\ref{teo:conta-iniettive}). Le funzioni iniettive
    da $I_k$ in $I_n$ sono $\frac{n!}{(n - k)!}$: sono i modi di \textbf{mettere in fila} $k$ elementi scelti tra $n$. Ma mettere in fila $k$
    elementi si fa in due passi:
    \begin{passi}
        \item si sceglie \textbf{quali} elementi: una $k$-parte, in $\binom{n}{k}$ modi;
        \item si sceglie \textbf{in che ordine} metterli: una permutazione dei $k$ elementi, in $k!$ modi (Corollario~\ref{cor:conta-biettive}).
    \end{passi}
    Quindi $\binom{n}{k} \cdot k! = \frac{n!}{(n - k)!}$, cioè $\binom{n}{k} = \frac{n!}{k!\,(n - k)!}$. Ad esempio, i podi (oro, argento, bronzo) con
    $10$ atleti sono $10 \cdot 9 \cdot 8 = 720$; ogni terna di atleti compare in $3! = 6$ podi diversi, quindi le terne sono $\binom{10}{3} = 720 / 6 = 120$.
\end{spiegazione}

\begin{esempio}{Contare sottoinsiemi}
    Sia $a = \{0, 1, \dots, 9\}$, con $10$ elementi.
    \begin{itemize}
        \item Le parti di $a$ con $3$ elementi sono $\binom{10}{3} = \frac{10!}{3!\,7!} = \frac{10 \cdot 9 \cdot 8}{6} = 120$; quelle con $11$ elementi sono $0$.
        \item Le parti di $a$ con $7$ elementi sono $\binom{10}{7} = \binom{10}{3} = 120$ (simmetria).
        \item Le parti formate da $0$ e da altri tre elementi distinti si ottengono scegliendo $3$ elementi tra i $9$ diversi da $0$:
              $\binom{9}{3} = \frac{9 \cdot 8 \cdot 7}{6} = 84$.
    \end{itemize}
\end{esempio}

% ──────────────────────────────────────────────────────────────
\section{Il teorema binomiale}\label{sec:newton}

I coefficienti binomiali devono il loro nome alla formula per la potenza di un \textbf{binomio} $a + b$.

\begin{concetto}
\begin{Teorema}{Teorema binomiale (formula di Newton)}{newton}
    \begin{prerequisiti}
        \dip{Teorema}{teo:stifel}
        \dip{Teorema}{teo:regole-potenze}
        \dip{Proposizione}{prop:multipli-unita}
        \dip{Corollario}{cor:induzione-uso}
    \end{prerequisiti}
    Sia $(s, +, \cdot)$ un anello unitario e siano $a, b \in s$ con $ab = ba$. Per ogni $n \in \mbN$:
    \[ (a + b)^n = \sum_{k = 0}^{n} \binom{n}{k} a^k b^{n - k} = b^n + \binom{n}{1} a b^{n - 1} + \dots + \binom{n}{n - 1} a^{n - 1} b + a^n. \]
\end{Teorema}
\begin{dimostrazione}
    \facoltativa
    Per induzione su $n$. \textbf{Base}: $(a + b)^0 = 1 = \binom{0}{0} a^0 b^0$. \textbf{Passo}: usando l'ipotesi induttiva, la distributività e
    $ab = ba$ (che permette di spostare le $a$ prima delle $b$),
    \[ (a + b)^{n + 1} = (a + b)\sum_{k = 0}^{n} \binom{n}{k} a^k b^{n - k} = \sum_{k = 0}^{n} \binom{n}{k} a^{k + 1} b^{n - k} + \sum_{k = 0}^{n} \binom{n}{k} a^k b^{n + 1 - k}. \]
    Nella prima somma il termine $a^j b^{n + 1 - j}$ compare con coefficiente $\binom{n}{j - 1}$, nella seconda con $\binom{n}{j}$: per la formula
    ricorsiva il coefficiente totale è $\binom{n}{j - 1} + \binom{n}{j} = \binom{n + 1}{j}$ (per $j = 0$ e $j = n + 1$ c'è un solo termine, con coefficiente
    $1 = \binom{n + 1}{0} = \binom{n + 1}{n + 1}$).
\end{dimostrazione}
\end{concetto}

\begin{spiegazione}
    Sviluppando $(a + b)^3 = (a + b)(a + b)(a + b)$ si sceglie, da ogni fattore, $a$ oppure $b$, e si moltiplicano le scelte. Il
    termine $a^2 b$ nasce da tutte le scelte con \textbf{due} $a$ e una $b$: tante quante le $2$-parti dei tre fattori, cioè $\binom{3}{2} = 3$
    ($aab$, $aba$, $baa$). Quindi
    \[ (a + b)^3 = b^3 + 3ab^2 + 3a^2 b + a^3, \]
    con i coefficienti $1, 3, 3, 1$ della riga $3$ del triangolo di Tartaglia. L'ipotesi $ab = ba$ serve a riordinare $aba$ e $baa$ in
    $a^2 b$: in un anello non commutativo la formula è falsa.
\end{spiegazione}
